\section{Yêu cầu bài toán}

\subsection{Miêu tả chi tiết bài toán}

Hệ thống phục vụ ba nhóm người dùng có nhu cầu rất khác nhau, và chính sự khác biệt đó tạo nên phần lớn độ phức tạp trong thiết kế và tối ưu.

Lễ tân và thu ngân làm việc trực tiếp tại quầy vào các khung giờ cao điểm nhận và trả phòng, đòi hỏi thao tác nhanh, quy trình tinh giản và phản hồi tức thì: màn hình phải trả lời ngay phòng nào trống, khách đã thanh toán chưa, số tiền còn thiếu là bao nhiêu. Độ trễ dù chỉ vài giây tại quầy cũng tác động trực tiếp đến sự chờ đợi của khách lưu trú.

Người quản lý và chủ cơ sở quan tâm đến các con số tổng hợp: doanh thu so với kỳ trước, công suất phòng và hiệu quả nguồn khách. Các truy vấn này quét trên tập dữ liệu lớn với yêu cầu khắt khe về tính nhất quán: số liệu giữa báo cáo phân tích và màn hình theo dõi thời gian thực cần đảm bảo tính đồng bộ, không để xảy ra sai lệch. Nhóm thứ ba gồm bộ phận kế toán và buồng phòng với các nghiệp vụ cụ thể như xuất hóa đơn, đối chiếu sổ quỹ, bàn giao ca và cập nhật tình trạng dọn phòng.

Qua quá trình làm việc và tiếp cận hệ thống, em nhận thấy sản phẩm cần đáp ứng một số yêu cầu chính như: số liệu cần đảm bảo tính nhất quán giữa các màn hình nghiệp vụ; quyền thao tác của người dùng phải được kiểm tra ở phía máy chủ, thay vì chỉ ẩn hoặc vô hiệu hóa nút trên giao diện; hệ thống cần hỗ trợ địa phương hóa tốt về ngôn ngữ, múi giờ và tiền tệ; đồng thời cần duy trì hiệu năng ổn định, tránh để các tác vụ xuất nhập dữ liệu lớn gây ảnh hưởng đến thao tác thường nhật của nhân viên tại quầy.

\begin{center}
\vspace{-4pt}
\begin{singlespace}
\begin{tikzpicture}[
  scale=0.88, every node/.append style={transform shape},
  box/.style={draw=xamvien!90!black, fill=xamnhat, rounded corners=3pt, align=center,
              font=\fontsize{8.5pt}{10pt}\selectfont, inner sep=2.5pt},
  nen/.style={draw=xanhkv, fill=white, rounded corners=3pt, align=center,
              font=\fontsize{8.5pt}{10pt}\selectfont, inner sep=2.5pt, thick},
  mui/.style={-{Stealth[length=2.5mm, width=1.6mm]}, draw=xanhkv!90!black, thick}
]

% ===== TẦNG 1: GIAO DIỆN NGƯỜI DÙNG (y = 0) =====
\node[box, text width=6.05cm, minimum height=0.72cm] (fe1) at (3.125, 0)
  {\textbf{\small Màn hình lễ tân}\\[1pt]{\scriptsize Thao tác tại quầy, phản hồi tức thì}};

\node[box, text width=9.25cm, minimum height=0.72cm] (fe2) at (11.25, 0)
  {\textbf{\small Màn hình quản lý}\\[1pt]{\scriptsize Theo dõi vận hành, giám sát và báo cáo}};

% ===== TẦNG 2: DỊCH VỤ NGHIỆP VỤ (y = -1.65) =====
\node[box, text width=2.85cm, minimum height=0.72cm] (sv1) at (1.50, -1.65)
  {\textbf{Dịch vụ đặt phòng}\\[1pt]{\scriptsize Xử lý đặt chỗ}};

\node[box, text width=2.85cm, minimum height=0.72cm] (sv2) at (4.75, -1.65)
  {\textbf{Dịch vụ quản lý}\\[1pt]{\scriptsize Danh mục, giá phòng}};

\node[box, text width=2.85cm, minimum height=0.72cm] (sv3) at (8.00, -1.65)
  {\textbf{Dịch vụ báo cáo}\\[1pt]{\scriptsize Thống kê số liệu}};

\node[box, text width=2.85cm, minimum height=0.72cm] (sv4) at (11.25, -1.65)
  {\textbf{Xuất nhập tệp}\\[1pt]{\scriptsize Xử lý tệp nền}};

\node[box, text width=2.85cm, minimum height=0.72cm] (sv5) at (14.50, -1.65)
  {\textbf{Giao tiếp di động}\\[1pt]{\scriptsize API cho ứng dụng}};

% ===== TẦNG 3: NỀN TẢNG DÙNG CHUNG (y = -3.20) =====
\node[nen, text width=15.70cm, minimum height=0.72cm] (nen) at (8.00, -3.20)
  {\textbf{\small Phần nền dùng chung}\\[1pt]{\scriptsize Thực thể nghiệp vụ, quy tắc phân quyền, cấu hình và thư viện tiện ích}};

% ===== TẦNG 4: CƠ SỞ DỮ LIỆU & LƯU TRỮ (y = -4.65) =====
\node[box, text width=7.65cm, minimum height=0.72cm] (db) at (3.925, -4.65)
  {\textbf{\small SQL Server}\\[1pt]{\scriptsize Lưu trữ có cấu trúc, phân vùng dữ liệu}};

\node[box, text width=7.65cm, minimum height=0.72cm] (redis) at (12.075, -4.65)
  {\textbf{\small Redis}\\[1pt]{\scriptsize Bộ nhớ đệm phân tán, điều phối khoá}};

% ===== MŨI TÊN TẦNG 1 -> TẦNG 2 (Rõ ràng, dài ~0.9cm) =====
\draw[mui] (sv1.north |- fe1.south) -- (sv1.north);
\draw[mui] ([xshift=-0.35cm]sv2.north |- fe1.south) -- ([xshift=-0.35cm]sv2.north);
\draw[mui] ([xshift=0.25cm]fe2.south west) -- ([xshift=0.35cm]sv2.north);
\draw[mui] (sv3.north |- fe2.south) -- (sv3.north);
\draw[mui] (sv4.north |- fe2.south) -- (sv4.north);
\draw[mui] (sv5.north |- fe2.south) -- (sv5.north);

% ===== MŨI TÊN TẦNG 2 -> TẦNG 3 (Dài ~0.85cm, thẳng đứng thanh thoát) =====
\draw[mui] (sv1.south) -- (nen.north -| sv1.south);
\draw[mui] (sv2.south) -- (nen.north -| sv2.south);
\draw[mui] (sv3.south) -- (nen.north -| sv3.south);
\draw[mui] (sv4.south) -- (nen.north -| sv4.south);
\draw[mui] (sv5.south) -- (nen.north -| sv5.south);

% ===== MŨI TÊN TẦNG 3 -> TẦNG 4 (Dài ~0.75cm, thẳng đứng dứt khoát) =====
\draw[mui] (nen.south -| db.north) -- (db.north);
\draw[mui] (nen.south -| redis.north) -- (redis.north);
\end{tikzpicture}
\end{singlespace}
\vspace{-4pt}
\captionof{figure}{Các thành phần chính của hệ thống và quan hệ giữa chúng}
\label{hinh:kientruc}
\vspace{-6pt}
\end{center}

\indent Hình \ref{hinh:kientruc} mô tả kiến trúc phân tầng và sự tương tác giữa các thành phần chính trong hệ thống. Trong quá trình thực tập, em nhận thấy các dịch vụ nghiệp vụ có mối liên kết chặt chẽ với nhau: một thao tác thực tế của người dùng thường đi qua nhiều tầng chức năng, do đó việc đọc hiểu mã nguồn hiện hữu là bước rất quan trọng trước khi bắt tay vào xây dựng tính năng mới.

\subsection{Phần bài toán em trực tiếp giải quyết}

Hệ thống sản phẩm là kết quả xây dựng của cả tập thể qua nhiều giai đoạn. Trong phạm vi thực tập, em được giao đảm nhiệm một số bài toán cụ thể với đầu vào và đầu ra rõ ràng, tập trung vào sáu hướng chính:

\begin{enumerate}[leftmargin=1.2cm, itemsep=2pt, topsep=3pt]
\item Tham gia xây dựng tính năng khóa phòng bảo trì: từ mô hình dữ liệu, quy tắc nghiệp vụ, API phía máy chủ đến giao diện người dùng, bản dịch đa ngôn ngữ và kiểm thử tự động.
\item Bổ sung hai chỉ tiêu thống kê về lượt khách nghỉ qua đêm và tổng số đêm nghỉ-bài toán xử lý khoảng thời gian và băm chuỗi định danh.
\item Tái cấu trúc phân quyền: tách quyền hủy đơn thành hai quyền độc lập, kèm kịch bản chuyển đổi dữ liệu phân quyền hai chiều cho các cơ sở đang hoạt động.
\item Tối ưu hiệu năng: cải tiến các câu truy vấn cơ sở dữ liệu, giải phóng khóa phân tán kịp thời và loại bỏ các phép tính toán lặp dư thừa trên giao diện.
\item Xây dựng các phân hệ cho bản quản lý mới: trang tổng quan (dashboard) với các biểu đồ doanh thu, danh mục phòng và tài liệu hướng dẫn tích hợp kênh bán phòng trực tuyến.
\item Xử lý đa ngôn ngữ và địa phương hóa: bổ sung tiếng Hàn, tiếng Trung, tiếng Lào cho các màn hình, đồng thời chuẩn hóa định dạng ngày tháng, tiền tệ và múi giờ.
\end{enumerate}

\subsection{Công việc của nhóm và công việc của em}

Nhóm phát triển vận hành theo mô hình đa chức năng: chuyên viên phân tích nghiệp vụ đặc tả yêu cầu; thiết kế trải nghiệm dựng giao diện; lập trình viên hiện thực logic; chuyên viên kiểm thử (QA) kiểm tra chất lượng trước phát hành; hỗ trợ kỹ thuật tiếp nhận phản hồi từ khách hàng; và trưởng nhóm điều phối, rà soát mã nguồn cùng tiến độ.

Trong cấu trúc đó, em tham gia với vai trò thực tập sinh phát triển phần mềm. Quy trình xử lý một phiếu công việc bắt đầu từ việc tiếp nhận yêu cầu, trao đổi làm rõ trường hợp biên, đọc mã nguồn hiện hữu để xác định phạm vi ảnh hưởng, viết mã, tự kiểm thử, bổ sung kiểm thử tự động cho logic quan trọng, gửi rà soát mã nguồn (code review) và phối hợp với bộ phận kiểm thử đến khi nghiệm thu hoàn tất.

Một bài học làm thay đổi nhận thức nghề nghiệp của em: nhiệm vụ thực sự của người phát triển không chỉ là làm cho tính năng chạy được, mà phải đảm bảo nó hoạt động chuẩn xác mà không gây ảnh hưởng tiêu cực đến bất kỳ hành vi nghiệp vụ nào đang vận hành ổn định trên hệ thống.
